\documentclass[12pt]{article}
\usepackage[utf8]{inputenc}
\usepackage[T1]{fontenc}
\usepackage{hyperref}
\title{Wasco-Wishram Language Resources for K-12 Curriculum Development}
\author{Compiled by Claude Sonnet 5 (claude-sonnet-5) for Andr\'as Kornai}
\date{}
\begin{document}
\maketitle

\section{Introduction}
Wasco-Wishram (ISO 639-3 wac; Glottocode wasc1239) is a Chinookan language, the variety of Upper Chinook (Kiksht) spoken at Warm Springs and Wishram. Sources disagree on its status. The Endangered Languages Project lists it as dormant, with a few elderly speakers reported in an apparently dated entry. Wikipedia reports the last fluent speakers as having died by 2006 and 2012.

\section{Primary Educational Resources}
Wikipedia says the University of Oregon Northwest Indian Language Institute partnered with Warm Springs schools to teach Kiksht. I did not verify this, and no actual teaching materials were found: \url{https://Www.wikipedia.org/wiki/Kiksht}

\section{Linguistic and Reference Materials}
\begin{itemize}
\item Endangered Languages Project entry: \url{https://endangeredlanguages.com/lang/1938}
\item Upper Chinook (Kiksht) article, which cites Dyk's \emph{A Grammar of Wishram} (1933) and Sapir and Curtin's \emph{Wishram Texts} (1909): \url{https://Www.wikipedia.org/wiki/Kiksht}
\end{itemize}
No current dictionary was found.

\section{Religious and Community Resources}
None found in the search results. The Confederated Tribes of Warm Springs would be the relevant community body. The search did not return its site.

\section{Digital and Audio Resources}
The Wikipedia article says audio and video recordings are held in the Endangered Languages Archive. I did not locate them.

\section{Curriculum Development Context}
This would be a heritage and revitalization curriculum built from archival sources (grammar, texts, recordings) with the Warm Springs and Yakama communities. Current speaker status and available materials need checking with them.
\end{document}